\documentclass[11pt,a4paper]{article}

% ---------- Кодировки и язык ----------
\usepackage[T2A]{fontenc}
\usepackage[utf8]{inputenc}
\usepackage[russian,english]{babel}

% ---------- Математика ----------
\usepackage{amsmath,amssymb,amsthm,mathtools}
\usepackage{bm}

% ---------- Оформление ----------
\usepackage[margin=2.5cm]{geometry}
\usepackage[protrusion=true,expansion=false]{microtype}
\usepackage{booktabs}
\usepackage{graphicx}
\usepackage{enumitem}
\usepackage{xcolor}
\usepackage{listings}
\usepackage{tikz}
\usepackage[colorlinks=true,linkcolor=blue!50!black,citecolor=green!50!black,urlcolor=blue!60!black]{hyperref}

% ---------- Теоремы ----------
\theoremstyle{plain}
\newtheorem{theorem}{Теорема}[section]
\newtheorem{proposition}[theorem]{Предложение}
\newtheorem{lemma}[theorem]{Лемма}
\newtheorem{hypothesis}{Гипотеза}
\theoremstyle{definition}
\newtheorem{definition}[theorem]{Определение}
\newtheorem{conjecture}[theorem]{Конъектура}
\theoremstyle{remark}
\newtheorem{remark}[theorem]{Замечание}

% ---------- Операторы и сокращения ----------
\DeclareMathOperator{\vol}{vol}
\DeclareMathOperator*{\argmin}{arg\,min}
\newcommand{\core}{\mathrm{core}}
\newcommand{\giant}{\mathrm{giant}}
\newcommand{\periph}{\mathrm{periph}}
\newcommand{\modu}{\mathrm{mod}}
\newcommand{\Ph}{\varphi}
\newcommand{\dmin}{d_{\min}}
\newcommand{\dmax}{d_{\max}}
\newcommand{\dbar}{\bar d}
\newcommand{\lamtwo}{\lambda_2}
\newcommand{\lamthree}{\lambda_3}
\newcommand{\Gcore}{G_{\core}}
\newcommand{\Ggiant}{G_{\giant}}
\newcommand{\Gtwo}{G_{2\core}}
\newcommand{\polylog}{\operatorname{polylog}}

% ---------- Листинги ----------
\lstset{
  basicstyle=\ttfamily\small,
  keywordstyle=\color{blue!60!black},
  commentstyle=\color{green!45!black},
  stringstyle=\color{red!60!black},
  showstringspaces=false,
  breaklines=true,
  frame=single,
  extendedchars=true,
  inputencoding=utf8,
  language=Python
}

% =====================================================================
\title{\bfseries Спектральная щель нормализованного лапласиана\\
в разреженных случайных графах с тяжёлохвостыми степенями:\\
асимптотика, bottleneck-структуры и связь с модульностью\\
биологических сетей}

\author{Кузьмин Леонид Леонидович\\[2pt]
\small магистр\\
\small\href{mailto:kuzminll@vk.com}{kuzminll@vk.com}}
\date{\today}

\begin{document}
\maketitle

\begin{abstract}
\noindent
Исследуется вторая собственная величина $\lamtwo$ нормализованного лапласиана
$L=I-D^{-1/2}AD^{-1/2}$ в разреженных случайных графах с тяжёлохвостым
распределением степеней $P(D=k)\sim Ck^{-\gamma}$ в режиме
$\dbar_{\min}=O(1)$, не покрытом классической теорией Chung--Lu--Vu.
Показывается, что одной зависимости от $\gamma$ недостаточно: порядок $\lamtwo$
определяется конкуренцией трёх масштабов --- низкостепенной периферии,
модульных разрезов и хабов/ядра. Вводится пара диагностических параметров
$(\Delta,\rho_*)$, где $\Delta=\Ph_{\modu}/\Ph_{\periph}$ --- отношение
модульного и периферийного conductance, а $\rho_*$ --- нормированный объём
минимального разреза. Предлагается фазовая диаграмма режимов малой
спектральной щели и численная схема (spectral sweep cut) для её построения.
Обсуждается осторожная биологическая интерпретация: малое $\lamtwo$ само по
себе не идентифицирует функциональную модульность.
\end{abstract}

\tableofcontents

% =====================================================================
\section{Мотивация и исходная постановка}
\label{sec:motivation}

\textbf{Исходная идея:} получить точную асимптотику второй собственной величины
$\lamtwo$ нормализованного лапласиана для разреженных случайных графов с
тяжёлохвостым распределением степеней $P(D=k)\sim Ck^{-\gamma}$ в режиме
$\dbar_{\min}=O(1)$, не покрытом теоремой Chung--Lu--Vu.

\textbf{Проблема исходной постановки.} Она слишком широкая и в части параметров
тривиально неверна. В разреженном режиме $\lamtwo$ может быть нулём из-за
несвязности, а после удаления мелких компонент --- иметь совершенно другой
порядок. Это не техническая мелочь, а центральная часть исследования.

\textbf{Корректировка.} Вместо одной формулы $\lamtwo\sim f(\gamma,\dbar,n)$
нужно строить \emph{фазовую диаграмму порядка $\lamtwo$} в зависимости от
\begin{equation}
\gamma,\qquad \dmin,\qquad \dmax,\qquad \dbar,\qquad n,
\end{equation}
и, что особенно важно, \emph{структуры низкостепенного хвоста}.

% =====================================================================
\section{Что уже известно из литературы}
\label{sec:known}

\begin{itemize}[leftmargin=*]
\item \textbf{Chung--Lu--Vu \cite{ChungLuVu2003}.} Для случайных графов с
заданными ожидаемыми степенями при $\dbar_{\min}\gg\ln^2 n$ спектр
нормализованного лапласиана концентрируется около предсказываемого поведения;
результаты включают power-law графы. Оценка щели:
\begin{equation}
\lambda(G)\;\ge\;1-4\,\dbar^{-1/2}-\dbar_{\min}^{-1}\ln^2 n .
\label{eq:CLV}
\end{equation}
При $\dbar_{\min}\le\ln^2 n$ оценка \eqref{eq:CLV} становится бессмысленной.

\item \textbf{Coja-Oghlan--Lanka \cite{CojaLanka2009}.} Для разреженного режима
$\dbar_{\min}=O(1)$ показано, что существует большой core, у которого
spectral gap нормализованного лапласиана остаётся отделённой от нуля:
\begin{equation}
\lamtwo(L_{\core})\;\ge\;1-c_0\,\dbar_{\min}^{-1/2}
\label{eq:COL}
\end{equation}
с высокой вероятностью. Речь идёт именно о \emph{core}, а не обо всём исходном
графе.

\item \textbf{Для $G(n,p)$ при ограниченной средней степени $d$.} Spectral gap
всего графа нормализованного лапласиана имеет $\lamtwo=o(1)$, хотя
giant/core может иметь gap порядка единицы.
\end{itemize}

\begin{remark}[ключевой вывод]
\begin{equation}
\boxed{\;\lamtwo(G)\;\not\sim\;\lamtwo(\core(G)).\;}
\end{equation}
\end{remark}

% =====================================================================
\section{Почему одного $\gamma$ недостаточно}
\label{sec:gamma}

Пусть $P(D=k)\sim Ck^{-\gamma}$. Это определяет хвост при больших $k$. Но
$\lamtwo$ очень чувствительна к \emph{маленьким степеням}.

\textbf{Контрпример.} Две последовательности с одинаковым $\gamma=2.5$:
\begin{itemize}[leftmargin=*]
\item модель A: $P(D=1)>0$;
\item модель B: $P(D=1)=0,\;P(D\ge 2)=1$.
\end{itemize}
У них одинаковый асимптотический хвост, но совершенно разное поведение
$\lamtwo$. В первой модели появляются листья и конечные компоненты; во второй
граф может оказаться существенно более связным.
\begin{equation}
\boxed{\;\gamma\ \text{не определяет}\ \lamtwo.\;}
\end{equation}

\textbf{Жёсткий контрпример.} Конфигурационная модель с $P(D=1)=p_1>0$. При
ограниченной средней степени такие графы естественно содержат мелкие
компоненты и dangling structures. Если граф несвязен, то $\lamtwo(L)=0$.
Следовательно, никакая нетривиальная формула вида $\lamtwo\sim f(\gamma)>0$
для всего графа невозможна.

% =====================================================================
\section{Три спектральные щели}
\label{sec:three}

\begin{table}[h]
\centering
\begin{tabular}{llp{7.2cm}}
\toprule
Объект & Обозначение & Что характеризует \\
\midrule
Полный граф & $\lamtwo(G)$ & Глобальная связность; $\lamtwo=0$ при несвязности \\
Giant component & $\lamtwo(\Ggiant)$ & Содержательная величина для sparse heavy-tailed сетей \\
2-core & $\lamtwo(\Gtwo)$ & Ненулевая asymptotic lower bound даже при $\dbar_{\min}=O(1)$ \\
\bottomrule
\end{tabular}
\caption{Три уровня спектральной щели.}
\end{table}

\begin{hypothesis}[порядок щелей]
\begin{equation}
\boxed{\;\lamtwo(G)\;\ll\;\lamtwo(\Ggiant)\;\lesssim\;\lamtwo(\Gtwo).\;}
\end{equation}
\end{hypothesis}

% =====================================================================
\section{Три механизма возникновения малой $\lamtwo$}
\label{sec:mechanisms}

\begin{enumerate}[leftmargin=*]
\item \textbf{Мелкие компоненты:} $\lamtwo(G)=0$.
\item \textbf{Dangling trees:} длинные «усы» создают малый conductance; по
Чигеру $\lamtwo\le 2\Ph$, поэтому достаточно длинный тонкий отросток может
сделать $\lamtwo\to 0$.
\item \textbf{Настоящая модульность:} две большие dense communities с малым
cut дают $\Ph\ll 1$ и $\lamtwo\ll 1$.
\end{enumerate}

\textbf{Ключевое разделение:}
\begin{equation}
\text{малое }\lamtwo\ \Longrightarrow\ \text{существует слабый глобальный cut},
\end{equation}
но
\begin{equation}
\text{малое }\lamtwo\ \not\Longrightarrow\ \text{биологическая модульность}.
\end{equation}
Нужно отделить \emph{community bottleneck} от \emph{peripheral bottleneck}.

% =====================================================================
\section{Формализация двух типов bottleneck}
\label{sec:bottleneck}

Пусть
\begin{equation}
\vol(S)=\sum_{v\in S}d_v,\qquad
\Ph(S)=\frac{|\partial S|}{\vol(S)} .
\end{equation}

\begin{definition}[периферийный bottleneck]
\begin{equation}
\Ph_{\periph}(G)=
\min_{\substack{S\subset V\\ 0<\vol(S)\le\varepsilon\,\vol(V)}}\Ph(S),
\end{equation}
где $\varepsilon\ll 1$, например $\varepsilon=0.05$.
\end{definition}

\begin{definition}[модульный bottleneck]
\begin{equation}
\Ph_{\modu}(G)=
\min_{\substack{S\subset V\\ \varepsilon\,\vol(V)\le\vol(S)\le\frac12\vol(V)}}\Ph(S).
\end{equation}
\end{definition}

Тогда
\begin{equation}
\Ph(G)=\min\{\Ph_{\periph},\Ph_{\modu}\}.
\end{equation}
Через неравенство Чигера \cite{Cheeger1970,KannanVempalaVetta2004}:
\begin{equation}
\frac{\Ph(G)^2}{2}\;\le\;\lamtwo\;\le\;2\,\Ph(G).
\label{eq:cheeger}
\end{equation}
То есть мы получаем не «разложение $\lamtwo$», а \emph{классификацию механизма,
создающего глобальный bottleneck}.

% =====================================================================
\section{Диагностические параметры}
\label{sec:diagnostics}

\begin{definition}[параметр конкуренции]
\begin{equation}
\boxed{\;\Delta(G)=\frac{\Ph_{\modu}(G)}{\Ph_{\periph}(G)}\;}
\end{equation}
\end{definition}
с интерпретацией:
\begin{itemize}[leftmargin=*]
\item $\Delta\gg 1$ --- периферийный bottleneck дешевле;
\item $\Delta\ll 1$ --- модульный bottleneck дешевле;
\item $\Delta\approx 1$ --- переходный режим.
\end{itemize}

Для устойчивости численных оценок использовать
\begin{equation}
\log\Delta=\log\Ph_{\modu}-\log\Ph_{\periph}.
\end{equation}

\begin{definition}[размер минимального спектрального cut]
\begin{equation}
\rho(S)=\frac{\vol(S)}{\vol(V)},\qquad
\rho_*=\rho(S_*),\qquad
S_*=\argmin_{S}\Ph(S).
\end{equation}
\end{definition}
Интерпретация:
\begin{itemize}[leftmargin=*]
\item $\rho_*=O(n^{-1})$ --- маленькая локальная структура;
\item $\rho_*=O(1)$ --- макроскопическое разделение графа.
\end{itemize}

\textbf{Диагностическая плоскость:} $(\log\Delta,\rho_*)$.

\textbf{Спектральный cut против комбинаторного.}
\begin{equation}
\rho_2=\frac{\vol(\{v:f_2(v)>0\})}{\vol(V)} .
\end{equation}
Если $\rho_2\approx\rho_*$ --- спектральный cut совпадает с комбинаторным.
Если нет --- $\lamtwo$ контролируется не минимальным cut'ом, а «усреднённым»
поведением.

\textbf{Множественные bottleneck'ы.} Разность $\lamthree-\lamtwo$: малая
разность означает несколько конкурирующих cut'ов с похожим conductance.

% =====================================================================
\section{Фазовая диаграмма (гипотетическая)}
\label{sec:phase}

\begin{figure}[h]
\centering
\begin{tikzpicture}[scale=1.15]
\draw[->,thick] (-0.2,0) -- (6.2,0) node[right] {$\log\Delta$};
\draw[->,thick] (0,-0.2) -- (0,4.2) node[above] {$\rho_*$};
\draw[dashed,gray] (0,2.1) -- (6.0,2.1);
\node[align=center] at (1.5,3.2) {\textsc{Modular}};
\node[align=center] at (4.4,3.2) {\textsc{Mixed}};
\node[align=center] at (1.5,0.9) {\textsc{Peripheral}};
\node[align=center] at (4.4,0.9) {\textsc{Core-dominated}};
\node[gray,anchor=north east] at (6.0,2.1) {\small макроскопический cut};
\node[gray,anchor=south east] at (6.0,2.1) {\small локальный cut};
\end{tikzpicture}
\caption{Гипотетическая фазовая диаграмма в координатах $(\log\Delta,\rho_*)$.}
\label{fig:phase}
\end{figure}

% =====================================================================
\section{Гипотезы}
\label{sec:hypotheses}

\begin{hypothesis}[H1]
Для sparse heavy-tailed configuration model с $P(D=k)\sim k^{-\gamma}$
асимптотика $\lamtwo(G)$ определяется не только $\gamma$, но конкуренцией трёх
масштабов:
\begin{equation}
\text{low-degree structures}\ \leftrightarrow\
\text{community cuts}\ \leftrightarrow\
\text{hub/core structure}.
\end{equation}
\end{hypothesis}

\begin{hypothesis}[H2]
Для полного графа при наличии положительной доли низкостепенных вершин
$\lamtwo(G)\to 0$ может быть тривиальным следствием периферии/несвязности.
\end{hypothesis}

\begin{hypothesis}[H3]
После перехода к giant component или 2-core $\lamtwo(\Gtwo)$ может оставаться
отделённой от нуля в режиме, где классическая теория Chung--Lu--Vu
\cite{ChungLuVu2003} не применима.
\end{hypothesis}

\begin{hypothesis}[H4]
Если дополнительно наложить planted modular structure, то появляется другой
масштаб:
\begin{equation}
\lamtwo\asymp\Ph_{\mathrm{module}}
\end{equation}
с точностью до квадратов/линейных факторов, задаваемых неравенством Чигера
\eqref{eq:cheeger}.
\end{hypothesis}

% =====================================================================
\section{Что будет новым результатом}
\label{sec:result}

\textbf{Теорема (целевая):}
\begin{equation}
\lamtwo(G_n)=\Theta\bigl(f(n,\gamma)\bigr)
\end{equation}
в определённом режиме. Отдельно:
\begin{equation}
\lamtwo(G_n^{2\core})=\Theta\bigl(g(\gamma,\dbar)\bigr)
\end{equation}
при других условиях. Если точную асимптотику получить не удастся,
\begin{equation}
c\,f(n,\gamma)\;\le\;\lamtwo\;\le\;C\,f(n,\gamma)
\end{equation}
--- уже полноценный результат.

% =====================================================================
\section{Математическая постановка}
\label{sec:setup}

Конфигурационная модель \cite{BenderCanfield1978,MolloyReed1995,MolloyReed1998}
с детерминированной последовательностью $d_1,\dots,d_n$:
\begin{equation}
\sum_i d_i\equiv 0 \pmod 2,\qquad
N_k(n)=\#\{i:d_i=k\}\sim c\,n\,k^{-\gamma}.
\end{equation}
Варьируем:
\begin{equation}
\dmin=1,2,3,\dots,\qquad
\dmax\sim n^{1/(\gamma-1)}
\end{equation}
либо controlled cutoff.

\textbf{Методы:} message-passing подходы для произвольных degree distributions;
spectral methods для конфигурационных моделей
\cite{DorogovtsevGoltsevMendesSamukhin2003,Soderberg2002,BordenaveLelargeMassoulie2018}.

% =====================================================================
\section{Экспериментальная часть}
\label{sec:experiments}

Для каждого $(n,\gamma,\dmin)$ генерируем $\approx 100$ реализаций при
$n=10^3,10^4,10^5,\dots$. Для каждой реализации считаем
\begin{align}
&\lamtwo(G),\quad \lamtwo(\Ggiant),\quad \lamtwo(\Gtwo),\\
&\Ph(G),\quad \Ph(\Ggiant),\quad \Ph(\Gtwo),\\
&\dmax,\quad \dbar,\quad N_1,\quad N_2 .
\end{align}

Scaling: строим $\log\lamtwo$ против $\log n$. Если $\lamtwo\sim n^{-\alpha}$,
то
\begin{equation}
\log\lamtwo=-\alpha\log n+C,
\end{equation}
и $\alpha$ оценивается статистически.

\begin{table}[h]
\centering
\begin{tabular}{lp{8.5cm}}
\toprule
Ансамбль & Что проверяем \\
\midrule
Power-law $\dmin=1$ & периферия \\
Power-law $\dmin=2$ & цепочки / 2-core \\
Power-law $\dmin=3$ & экспандероподобный core \\
Power-law $+$ planted communities & настоящая модульность \\
\bottomrule
\end{tabular}
\caption{Четыре экспериментальных ансамбля.}
\end{table}

% =====================================================================
\section{Алгоритм вычисления $\Ph_{\periph}$ и $\Ph_{\modu}$}
\label{sec:algorithm}

\subsection*{Подход 1: спектральная аппроксимация через sweep cut}

\begin{enumerate}[leftmargin=*]
\item Вычислить второй собственный вектор $f_2$ нормализованного лапласиана
$L=I-D^{-1/2}AD^{-1/2}$.
\item Отсортировать вершины по $f_2(v)$: $v_1,\dots,v_n$.
\item Sweep cut: для каждого $k=1,\dots,n-1$ рассмотреть
$S_k=\{v_1,\dots,v_k\}$ и вычислить $\Ph(S_k)$.
\item Найти минимумы:
\begin{align}
\Ph_{\periph}^{\mathrm{spec}}
&=\min\{\Ph(S_k):\vol(S_k)\le\varepsilon\vol(V)\},\\
\Ph_{\modu}^{\mathrm{spec}}
&=\min\{\Ph(S_k):\varepsilon\vol(V)\le\vol(S_k)\le\tfrac12\vol(V)\}.
\end{align}
\item Записать соответствующие $\rho_*$.
\end{enumerate}

\textbf{Сложность:} $O(m\cdot\mathrm{iter}+n\log n)$.

\subsection*{Подход 2: local Cheeger с seed-множествами}

\begin{enumerate}[leftmargin=*]
\item Выбрать seed-множества (низкостепенные вершины для периферии; хабы или
случайные вершины для модульности).
\item Для каждого seed применить local clustering (Andersen--Chung--Lang,
PageRank-based) \cite{AndersenChungLang2006}.
\item Собрать все найденные множества $S$, вычислить $\Ph(S)$.
\item Отфильтровать по объёму.
\end{enumerate}

\textbf{Сложность:} $O(nm\log n/\Ph^2)$.

\textbf{Рекомендация:} начать со спектрального sweep cut. Если
$\rho_2\ne\rho_*$, добавить local clustering для проверки.

% =====================================================================
\section{Псевдокод sweep cut}
\label{sec:pseudocode}

\begin{lstlisting}
import numpy as np
from scipy.sparse.linalg import eigsh

def spectral_diagnostics(L, degrees, adj, eps=0.05):
    n = len(degrees)
    vol_V = degrees.sum()

    vals, vecs = eigsh(L, k=2, which='SM')
    lambda2 = vals[1]
    f2 = vecs[:, 1]

    order = np.argsort(f2)

    vol_S = 0
    cut_size = 0
    in_S = np.zeros(n, dtype=bool)

    phi_periph = np.inf
    phi_mod = np.inf
    rho_periph = None
    rho_mod = None

    for k in range(n - 1):
        v = order[k]
        in_S[v] = True
        vol_S += degrees[v]

        for u in adj[v]:
            if in_S[u]:
                cut_size -= 1
            else:
                cut_size += 1

        if vol_S == 0:
            continue
        phi = cut_size / vol_S

        if vol_S <= eps * vol_V:
            if phi < phi_periph:
                phi_periph = phi
                rho_periph = vol_S / vol_V
        elif vol_S <= 0.5 * vol_V:
            if phi < phi_mod:
                phi_mod = phi
                rho_mod = vol_S / vol_V

    log_Delta = np.log(phi_mod) - np.log(phi_periph) if phi_periph > 0 else np.inf

    return {'lambda2': lambda2, 'phi_periph': phi_periph,
            'phi_mod': phi_mod, 'rho_periph': rho_periph,
            'rho_mod': rho_mod, 'log_Delta': log_Delta}
\end{lstlisting}

Реализация в репозитории: \texttt{src/diagnostics.py}.

% =====================================================================
\section{Биологическая интерпретация}
\label{sec:bio}

\textbf{Осторожность со словом «устойчивость».} $\lamtwo$ характеризует
глобальную связность/expansion и через Чигера связан с bottleneck'ами. Для
random walk $P=D^{-1}A$ и $\mu_2=1-\lamtwo$. Малый $\lamtwo$ соответствует
медленной релаксации случайного блуждания. Это можно интерпретировать как
\emph{структурную интегрированность} сети, но не автоматически как
биологическую robustness.

\textbf{Нужно определить, какую именно устойчивость изучаем:}
\begin{itemize}[leftmargin=*]
\item устойчивость к удалению узлов;
\item устойчивость к удалению рёбер;
\item сохранение giant component;
\item устойчивость потоков/диффузии;
\item устойчивость функциональных модулей.
\end{itemize}

\textbf{Масштаб функционального разделения через $\rho_*$:}
\begin{itemize}[leftmargin=*]
\item $\rho_*\approx 0.5$ --- две большие половины (два полушария, два
метаболических режима);
\item $\rho_*\approx 0.1$ --- маленький модуль, отделённый от остального;
\item $\rho_*\approx n^{-1}$ --- один узел или маленькая группа.
\end{itemize}

\begin{quote}
\itshape
В биологических сетях малая $\lamtwo$ сама по себе не идентифицирует
функциональную модульность. Предлагаемый анализ дополнительно определяет
масштаб минимального bottleneck и позволяет различать локальную периферию и
макроскопическое разделение сети.
\end{quote}

% =====================================================================
\section{Связь с хабами}
\label{sec:hubs}

Тяжёлый хвост создаёт $\dmax\gg\dbar$. Хабы сильно влияют на спектр adjacency
matrix: необычно высокие степени могут создавать изолированные собственные
значения с локализованными собственными векторами около хабов. Но это
\emph{не означает}, что хабы обязательно увеличивают $\lamtwo$.

\textbf{Важное различие:}
\begin{equation}
\lambda_{\max}(A)\qquad\text{и}\qquad\lamtwo(L)
\end{equation}
чувствуют разные свойства сети. Хаб может резко изменить верхнюю часть
спектра, практически не решив глобальный bottleneck.

% =====================================================================
\section{Четырёхуровневая структура проекта}
\label{sec:levels}

\begin{description}[leftmargin=*]
\item[Уровень I --- математическая модель.] Configuration model:
$(d_1,\dots,d_n)$, $N_k(n)\sim c\,n\,k^{-\gamma}$. Варьируем
$\gamma,\dmin,\dmax,n$.
\item[Уровень II --- спектральная асимптотика.] Исследуем $\lamtwo(G)$,
$\lamtwo(\Ggiant)$, $\lamtwo(\Gtwo)$.
\item[Уровень III --- механизм возникновения щели.] Измеряем
$\Ph_{\periph},\Ph_{\modu},\Delta,\rho_*$. Проверяем
$\lamtwo\leftrightarrow\min(\Ph_{\periph},\Ph_{\modu})$.
\item[Уровень IV --- биологическая интерпретация.] Проверяем, соответствует ли
наблюдаемая малая щель peripheral sparsity или macroscopic modular
organization. И только после этого обсуждаем implications для
robustness / functional integration.
\end{description}

% =====================================================================
\section{Главный исследовательский вопрос}
\label{sec:question}

\begin{quote}
Как низкостепенная периферия и тяжёлый хвост распределения степеней совместно
определяют асимптотику второй собственной величины нормализованного
лапласиана в разреженных случайных графах, и как с помощью пары
$(\Delta,\rho_*)$ различить уменьшение $\lamtwo$, вызванное периферийным
bottleneck, от уменьшения, вызванного макроскопической модульностью?
\end{quote}

Потенциальный результат --- не просто формула, а \emph{теория режимов
возникновения малой спектральной щели}.

% =====================================================================
\section{Открытые вопросы и риски}
\label{sec:open}

\begin{itemize}[leftmargin=*]
\item \textbf{Вычислимость $\Ph$:} точное вычисление NP-трудно; нужна явная
аппроксимация (sweep cut или local clustering).
\item \textbf{Устойчивость $\Delta$:} использовать $\log\Delta$.
\item \textbf{Связь с собственным вектором:} сравнить $\rho_2$ и $\rho_*$.
\item \textbf{Множественные bottleneck'ы:} добавить $\lamthree-\lamtwo$.
\item \textbf{2-core против giant component:} уточнить, когда совпадают, когда
различаются.
\item \textbf{Размазанная периферия:} если low-degree узлы не компактны,
$\Ph_{\periph}$ может быть неинформативен; проверить численно.
\item \textbf{Аналитическая часть:} может оказаться неподъёмной; план B ---
численные гипотезы $+$ частичные оценки.
\item \textbf{Биологическая интерпретация:} требует осторожности; не называть
$\lamtwo$ мерой «устойчивости» без уточнения.
\end{itemize}

% =====================================================================
\section{План дальнейших действий}
\label{sec:plan}

\begin{enumerate}[leftmargin=*]
\item Реализовать генерацию графов для четырёх ансамблей.
\item Реализовать sweep cut для $\Ph_{\periph},\Ph_{\modu},\Delta,\rho_*$.
\item Провести численные эксперименты для $n=10^3,10^4,10^5$.
\item Построить фазовую диаграмму в $(\log\Delta,\rho_*)$.
\item Проверить гипотезу
$\lamtwo\leftrightarrow\min(\Ph_{\periph},\Ph_{\modu})$.
\item Сформулировать эмпирические гипотезы о порядке $\lamtwo$ в каждом режиме.
\item Попробовать аналитический вывод для простейшего случая
(смешанно-регулярный граф, два типа степеней).
\item Оформить результаты как препринт (arXiv: math.SP, math.CO, q-bio.MN).
\end{enumerate}

% =====================================================================
\appendix
\section{Численная реализация}
\label{app:code}

Код проекта:
\begin{center}
\begin{tabular}{ll}
\toprule
Файл & Назначение \\
\midrule
\texttt{run\_experiments.py} & оркестрация 4 ансамблей, CSV-вывод, оценка $\alpha$ \\
\texttt{src/generators.py} & степенные последовательности, erased configuration model, planted block model \\
\texttt{src/diagnostics.py} & $\lamtwo$ (G/giant/2-core), sweep cut, $\Ph,\Delta,\rho$ \\
\bottomrule
\end{tabular}
\end{center}

Малые собственные значения вычисляются устойчиво через
$M=I-\tfrac12 L=\tfrac12\,(I+D^{-1/2}AD^{-1/2})$: малые $\lambda$ соответствуют
большим собственным значениям $M$, поэтому используется
\texttt{eigsh(M, which="LA")} с последующим преобразованием
$\lambda=2(1-\mu)$.

% =====================================================================
\begin{thebibliography}{99}

\bibitem{BenderCanfield1978}
E.~A.~Bender, E.~R.~Canfield.
\newblock The asymptotic number of labeled graphs with given degree sequences.
\newblock \emph{Journal of Combinatorial Theory, Series A} \textbf{24}(3):296--307, 1978.

\bibitem{BollobasRiordan2004}
B.~Bollob\'as, O.~Riordan.
\newblock The diameter of a scale-free random graph.
\newblock \emph{Combinatorica} \textbf{24}(1):5--34, 2004.

\bibitem{BordenaveLelargeMassoulie2018}
C.~Bordenave, M.~Lelarge, L.~Massouli\'e.
\newblock Non-backtracking spectrum of random graphs: community detection and
non-regular Ramanujan graphs.
\newblock \emph{Annals of Probability} \textbf{46}(1):1--71, 2018.

\bibitem{Chung1997}
F.~Chung.
\newblock \emph{Spectral Graph Theory}.
\newblock CBMS Regional Conference Series in Mathematics, vol.~92, AMS, 1997.

\bibitem{ChungLu2006}
F.~Chung, L.~Lu.
\newblock \emph{Complex Graphs and Networks}.
\newblock CBMS Regional Conference Series in Mathematics, vol.~107, AMS, 2006.

\bibitem{ChungLuVu2003}
F.~Chung, L.~Lu, V.~Vu.
\newblock Spectra of random graphs with given expected degrees.
\newblock \emph{Proceedings of the National Academy of Sciences}
\textbf{100}(11):6313--6318, 2003.

\bibitem{ChungRadcliffe2011}
F.~Chung, M.~Radcliffe.
\newblock On the spectra of general random graphs.
\newblock \emph{Electronic Journal of Combinatorics} \textbf{18}(1), P215, 2011.

\bibitem{CojaLanka2009}
A.~Coja-Oghlan, S.~Lanka.
\newblock Finding planted partitions in random graphs with general degree
distributions.
\newblock \emph{SIAM Journal on Discrete Mathematics} \textbf{23}(4):1682--1714, 2009.

\bibitem{CojaLanka2009b}
A.~Coja-Oghlan, S.~Lanka.
\newblock The spectral gap of random graphs with given expected degrees.
\newblock \emph{Electronic Journal of Combinatorics} \textbf{16}(1), R138, 2009.

\bibitem{Cheeger1970}
J.~Cheeger.
\newblock A lower bound for the smallest eigenvalue of the Laplacian.
\newblock In \emph{Problems in Analysis}, Princeton University Press, 195--199, 1970.

\bibitem{KannanVempalaVetta2004}
R.~Kannan, S.~Vempala, A.~Vetta.
\newblock On clusterings: good, bad and spectral.
\newblock \emph{Journal of the ACM} \textbf{51}(3):497--515, 2004.

\bibitem{AroraRaoVazirani2009}
S.~Arora, S.~Rao, U.~Vazirani.
\newblock Expander flows, geometric embeddings and graph partitioning.
\newblock \emph{Journal of the ACM} \textbf{56}(2), 2009.

\bibitem{AndersenChungLang2006}
R.~Andersen, F.~Chung, K.~Lang.
\newblock Local graph partitioning using PageRank vectors.
\newblock \emph{FOCS 2006}, 475--486.

\bibitem{LeeOveisGharanTrevisan2012}
J.~R.~Lee, S.~Oveis Gharan, L.~Trevisan.
\newblock Multi-way spectral partitioning and higher-order Cheeger inequalities.
\newblock \emph{STOC 2012}, 1117--1130.

\bibitem{MolloyReed1995}
M.~Molloy, B.~Reed.
\newblock A critical point for random graphs with a given degree sequence.
\newblock \emph{Random Structures \& Algorithms} \textbf{6}(2--3):161--180, 1995.

\bibitem{MolloyReed1998}
M.~Molloy, B.~Reed.
\newblock The size of the giant component of a random graph with a given degree
sequence.
\newblock \emph{Combinatorics, Probability and Computing} \textbf{7}(3):295--305, 1998.

\bibitem{PittelSpencerWormald1996}
B.~Pittel, J.~Spencer, N.~Wormald.
\newblock Sudden emergence of a giant $k$-core in a random graph.
\newblock \emph{Journal of Combinatorial Theory, Series B} \textbf{67}(1):111--151, 1996.

\bibitem{JansonLuczakRucinski2000}
S.~Janson, T.~{\L}uczak, A.~Ruci\'nski.
\newblock \emph{Random Graphs}.
\newblock Wiley-Interscience, 2000.

\bibitem{BarabasiAlbert1999}
A.-L.~Barab\'asi, R.~Albert.
\newblock Emergence of scaling in random networks.
\newblock \emph{Science} \textbf{286}(5439):509--512, 1999.

\bibitem{Faloutsos1999}
M.~Faloutsos, P.~Faloutsos, C.~Faloutsos.
\newblock On power-law relationships of the Internet topology.
\newblock \emph{SIGCOMM 1999}, 251--262.

\bibitem{DorogovtsevGoltsevMendesSamukhin2003}
S.~N.~Dorogovtsev, A.~V.~Goltsev, J.~F.~F.~Mendes, A.~N.~Samukhin.
\newblock Spectra of complex networks.
\newblock \emph{Physical Review E} \textbf{68}, 046109, 2003.

\bibitem{Soderberg2002}
B.~S\"oderberg.
\newblock General formalism for inhomogeneous random graphs.
\newblock \emph{Physical Review E} \textbf{66}, 066121, 2002.

\bibitem{NewmanGirvan2004}
M.~E.~J.~Newman, M.~Girvan.
\newblock Finding and evaluating community structure in networks.
\newblock \emph{Physical Review E} \textbf{69}, 026113, 2004.

\bibitem{Newman2006}
M.~E.~J.~Newman.
\newblock Modularity and community structure in networks.
\newblock \emph{PNAS} \textbf{103}(23):8577--8582, 2006.

\bibitem{DecelleKrzakalaMooreZdeborova2011}
A.~Decelle, F.~Krzakala, C.~Moore, L.~Zdeborov\'a.
\newblock Asymptotic analysis of the stochastic block model for modular networks
and its algorithmic applications.
\newblock \emph{Physical Review E} \textbf{84}, 066106, 2011.

\bibitem{Abbe2018}
E.~Abbe.
\newblock Community detection and stochastic block models: recent developments.
\newblock \emph{Journal of Machine Learning Research} \textbf{18}(177):1--86, 2018.

\bibitem{AlbertJeongBarabasi2000}
R.~Albert, H.~Jeong, A.-L.~Barab\'asi.
\newblock Error and attack tolerance of complex networks.
\newblock \emph{Nature} \textbf{406}:378--382, 2000.

\bibitem{AlbertBarabasi2002}
R.~Albert, A.-L.~Barab\'asi.
\newblock Statistical mechanics of complex networks.
\newblock \emph{Reviews of Modern Physics} \textbf{74}(1):47--97, 2002.

\bibitem{Newman2010}
M.~E.~J.~Newman.
\newblock \emph{Networks: An Introduction}.
\newblock Oxford University Press, 2010.

\bibitem{BarratBarthelemyVespignani2008}
A.~Barrat, M.~Barth\'elemy, A.~Vespignani.
\newblock \emph{Dynamical Processes on Complex Networks}.
\newblock Cambridge University Press, 2008.

\end{thebibliography}

\end{document}
